\documentclass[12pt,a4paper]{article}

\usepackage[T2A]{fontenc}
\usepackage[utf8]{inputenc}
\usepackage[russian]{babel}
\usepackage{amsmath,amssymb}
\usepackage{graphicx}
\usepackage{tikz}
\usepackage{xcolor}
\usepackage{hyperref}
\usepackage{geometry}

\geometry{left=25mm,right=20mm,top=20mm,bottom=20mm}
\allowdisplaybreaks

\newcommand{\todo}[1]{%
    \par\noindent\textcolor{red}{\textbf{TODO:} #1}\par
}
\newcommand{\Rbase}{R_0}
\newcommand{\Uref}{U_0}
\newcommand{\rhoref}{\rho_0}

\title{Параметрическая физически-информированная нейронная сеть для задачи всплытия пузырька}
\author{Клявлин Эмиль Рамилевич}
\date{2026}

\begin{document}

\maketitle

\section{Введение}

Моделирование двухфазных течений с подвижной границей раздела является важной
задачей вычислительной гидродинамики. Такие процессы встречаются при движении
капель и пузырьков, в теплообменных аппаратах, химических реакторах и других
инженерных системах. Численное решение уравнений двухфазного течения обычно
требует мелких сеток около интерфейса и значительных вычислительных ресурсов.
Поэтому представляет интерес построение моделей, которые используют физические
законы и ограниченный набор данных CFD-расчетов для восстановления полей
скорости, давления и положения интерфейса.

В данной работе рассматривается тестовая задача \textit{rising bubble}: легкий
пузырек находится внутри более тяжелой жидкости и всплывает под действием силы
тяжести. Базовая постановка взята из работы Buhendwa, Adami и Adams
\cite{Buhendwa2021}, где физически-информированные нейронные сети применялись
для восстановления полей несжимаемого двухфазного течения по данным о движении
интерфейса. В отличие от исходной однорадиусной постановки, в настоящей работе
строится параметрическая модель: радиус начального пузырька \(R\) является
дополнительным входом нейронной сети. Цель такой модели --- предсказывать
динамику всплытия для радиусов, которые не использовались при обучении.

\subsection{Цель и задачи}

Цель работы состоит в разработке и исследовании параметрической
физически-информированной нейронной сети для восстановления двумерного
нестационарного течения в задаче всплытия пузырька при изменяющемся начальном
радиусе \(R\).

Для достижения цели решаются следующие задачи:
\begin{enumerate}
    \item сгенерировать CFD-данные для нескольких радиусов начального пузырька
    с помощью Basilisk;
    \item привести данные к единому HDF5-формату, содержащему сетку, время,
    level-set, плотность, давление и компоненты скорости;
    \item построить PINN-модель, которая принимает на вход координаты, время и
    параметр радиуса;
    \item обучить модель на наборе радиусов и проверить качество обобщения на
    радиусах, исключенных из обучающей выборки;
    \item сравнить предсказания PINN с CFD-решением по положению интерфейса и
    полям \(u\), \(v\), \(p\).
\end{enumerate}

\section{Математическая постановка}

\subsection{Геометрия и физические параметры}

Рассматривается двумерная область
\[
    (x,y)\in[-0.5,0.5]\times[0,2],
    \qquad
    t\in[0,3].
\]
Эта область эквивалентна области
\([-0.5,0.5]\times[-1,1]\), использованной в статье \cite{Buhendwa2021},
с точностью до вертикального сдвига. Начальный центр пузырька в используемых
данных находится в точке
\[
    (x_c,y_c)=(0,0.5).
\]
Начальный радиус \(R\) меняется от расчета к расчету. В проекте используются
следующие наборы:
\[
    R_{\mathrm{train}}\in\{0.18,0.22,0.25,0.28,0.32\},
    \qquad
    R_{\mathrm{test}}\in\{0.20,0.30\}.
\]
Таким образом, радиусы \(0.20\) и \(0.30\) являются невидимыми для модели во
время обучения и используются для проверки параметрического обобщения.

Плотности и динамические вязкости фаз совпадают с параметрами rising-bubble
теста из \cite{Buhendwa2021}:
\[
    \rho_1=100,\qquad \rho_2=1000,
    \qquad
    \mu_1=1,\qquad \mu_2=10.
\]
Здесь фаза \(1\) соответствует легкому пузырьку, а фаза \(2\) --- окружающей
жидкости. Коэффициент поверхностного натяжения и ускорение свободного падения
равны
\[
    \sigma=24.5,
    \qquad
    \mathbf{g}=(0,-0.98)^T.
\]

На верхней и нижней границах задается условие прилипания:
\[
    u=0,\qquad v=0.
\]
На боковых границах в PINN-постановке используется периодическое условие:
\[
    u(-0.5,y,t)=u(0.5,y,t),\quad
    v(-0.5,y,t)=v(0.5,y,t),\quad
    p(-0.5,y,t)=p(0.5,y,t).
\]
На верхней границе дополнительно фиксируется безразмерное опорное давление
\[
    p_\infty^*=1.
\]
Это необходимо, поскольку в несжимаемой постановке давление определяется с
точностью до аддитивной константы.

\begin{figure}[h]
    \centering
    \begin{tikzpicture}[scale=0.85,>=stealth]
        \fill[gray!20] (0,0) rectangle (4,8);
        \draw[very thick] (0,0) rectangle (4,8);
        \fill[white] (2,2) circle (0.95);
        \draw[thick] (2,2) circle (0.95);
        \node at (2,2) {fluid 1};
        \node at (1.05,0.55) {fluid 2};
        \node at (2,8.35) {$u=v=0,\ p_\infty^*=1$};
        \node at (2,-0.35) {$u=v=0$};
        \node[rotate=90] at (-0.35,4) {periodic};
        \node[rotate=90] at (4.35,4) {periodic};
        \draw[->,thick] (3.15,6.5) -- (3.15,5.35);
        \node[right] at (3.15,5.9) {$\mathbf{g}$};
        \draw[<->] (0,8.8) -- (4,8.8);
        \node[fill=white,inner sep=2pt] at (2,8.8) {$1.0$};
        \draw[<->] (-0.9,0) -- (-0.9,8);
        \node[fill=white,inner sep=2pt] at (-0.9,4) {$2.0$};
        \draw[<->] (2.95,2) -- (2,2);
        \node[fill=white,inner sep=2pt] at (2.48,2.25) {$R$};
        \draw[dashed] (2,0) -- (2,2);
        \node[right] at (2.05,0.9) {$0.5$};
    \end{tikzpicture}
    \caption{Схема начального положения пузырька и граничных условий.}
    \label{fig:setup}
\end{figure}

\subsection{Безразмерные переменные}

В исходной статье для одного радиуса использовался масштаб длины \(L_r=R=0.25\).
В параметрической задаче такой выбор сделал бы безразмерный радиус одинаковым
для всех расчетов. Поэтому в данной работе используется фиксированный базовый
масштаб
\[
    \Rbase=0.25,\qquad \Uref=1,\qquad \rhoref=\rho_2=1000.
\]
Безразмерные переменные определяются как
\[
    x^*=\frac{x}{\Rbase},\qquad
    y^*=\frac{y}{\Rbase},\qquad
    t^*=\frac{t\Uref}{\Rbase},\qquad
    p^*=\frac{p}{\rhoref \Uref^2}.
\]
Радиус подается в нейронную сеть как дополнительный безразмерный параметр
\[
    R^*=\frac{R}{\Rbase}.
\]
Для обучающих и тестовых радиусов это дает
\[
    R^*_{\mathrm{train}}\in\{0.72,0.88,1.00,1.12,1.28\},
    \qquad
    R^*_{\mathrm{test}}\in\{0.80,1.20\}.
\]

\subsection{Уравнения двухфазного течения}

В качестве фазовой переменной используется объемная доля легкой фазы
\(\alpha(x,y,t)\). Для пузырька \(\alpha=1\), для окружающей жидкости
\(\alpha=0\), а область интерфейса соответствует \(0<\alpha<1\).
Плотность и вязкость задаются через \(\alpha\):
\begin{align}
    \rho(\alpha) &= \rho_2 + (\rho_1-\rho_2)\alpha, \\
    \mu(\alpha) &= \mu_2 + (\mu_1-\mu_2)\alpha.
\end{align}

Безразмерная форма уравнений совпадает с постановкой \cite{Buhendwa2021}.
Условие несжимаемости:
\begin{equation}
    \nabla^*\cdot\mathbf{u}^*=0.
    \label{eq:mass}
\end{equation}
Уравнение переноса объемной доли:
\begin{equation}
    \frac{\partial\alpha}{\partial t^*}
    + \mathbf{u}^*\cdot\nabla^*\alpha=0.
    \label{eq:alpha_transport}
\end{equation}
Уравнение импульса с поверхностным натяжением и силой тяжести:
\begin{align}
    \rho^*
    \left(
        \frac{\partial\mathbf{u}^*}{\partial t^*}
        +(\mathbf{u}^*\cdot\nabla^*)\mathbf{u}^*
    \right)
    &=
    -\nabla^*p^*
    +\nabla^*\cdot
    \left[
        C_\mu(\alpha)
        \left(\nabla^*\mathbf{u}^*+\nabla^*\mathbf{u}^{*T}\right)
    \right]
    \notag\\
    &\quad
    + C_\sigma \kappa\nabla^*\alpha
    + \rho^*\mathbf{C}_g.
    \label{eq:momentum}
\end{align}
Здесь
\[
    \rho^*=\frac{\rho}{\rhoref},
    \qquad
    C_\mu(\alpha)=\frac{\mu(\alpha)}{\rhoref\Uref\Rbase},
    \qquad
    C_\sigma=\frac{\sigma}{\rhoref\Uref^2\Rbase},
\]
\[
    \mathbf{C}_g=\frac{\mathbf{g}\Rbase}{\Uref^2}.
\]
При выбранных параметрах \(C_\sigma=0.098\), что соответствует
\(We=1/C_\sigma\approx 10.2\), как в статье \cite{Buhendwa2021}.

Кривизна интерфейса вычисляется по полю \(\alpha\):
\begin{equation}
    \kappa
    =
    -\nabla^*\cdot
    \frac{\nabla^*\alpha}{|\nabla^*\alpha|}.
    \label{eq:curvature_vector}
\end{equation}
В реализации используется развернутая форма с малым \(\varepsilon\), защищающим
от деления на ноль:
\begin{align}
    \kappa
    =
    -\left[
    \frac{\alpha_{xx}+\alpha_{yy}}{\sqrt{\alpha_x^2+\alpha_y^2+\varepsilon}}
    -
    \frac{
        \alpha_x^2\alpha_{xx}
        +2\alpha_x\alpha_y\alpha_{xy}
        +\alpha_y^2\alpha_{yy}
    }{
        \left(\alpha_x^2+\alpha_y^2+\varepsilon\right)^{3/2}
    }
    \right].
    \label{eq:curvature_expanded}
\end{align}
Все производные в формулах ниже берутся по безразмерным входам
\((x^*,y^*,t^*)\).

\subsection{Физические невязки}

PINN обучается не напрямую решать дискретную CFD-задачу, а минимизировать
невязки уравнений \eqref{eq:mass}--\eqref{eq:momentum}. Для компактности далее
звездочка у безразмерных переменных опускается. Невязка несжимаемости:
\begin{equation}
    r_m=u_x+v_y.
\end{equation}
Невязка переноса объемной доли:
\begin{equation}
    r_\alpha=\alpha_t+u\alpha_x+v\alpha_y.
\end{equation}

В коде переменная вязкость раскрывается покомпонентно; при записи ниже
используется условие несжимаемости \(u_x+v_y=0\). Обозначим
\[
    C_{\mu,x}=\frac{\mu_x}{\rhoref\Uref\Rbase},
    \qquad
    C_{\mu,y}=\frac{\mu_y}{\rhoref\Uref\Rbase},
    \qquad
    C_{g,y}=\frac{g_y\Rbase}{\Uref^2}.
\]
Тогда невязки уравнений импульса имеют вид:
\begin{align}
    r_x
    &=
    \rho^*
    \left(u_t+uu_x+vu_y\right)
    +p_x
    -C_\sigma\kappa\alpha_x
    -C_\mu\left(u_{xx}+u_{yy}\right)
    \notag\\
    &\quad
    -2C_{\mu,x}u_x
    -C_{\mu,y}\left(u_y+v_x\right),
    \label{eq:rx}
    \\[0.4em]
    r_y
    &=
    \rho^*
    \left(v_t+uv_x+vv_y\right)
    +p_y
    -C_\sigma\kappa\alpha_y
    -C_\mu\left(v_{xx}+v_{yy}\right)
    \notag\\
    &\quad
    -\rho^*C_{g,y}
    -2C_{\mu,y}v_y
    -C_{\mu,x}\left(u_y+v_x\right).
    \label{eq:ry}
\end{align}
Знак гравитационного члена в \eqref{eq:ry} соответствует переносу правой части
уравнения импульса в левую часть. Так как \(g_y=-0.98\), коэффициент
\(C_{g,y}\) отрицателен.

\section{Данные CFD}

\subsection{Генерация данных в Basilisk}

Для генерации данных используется файл
\texttt{diploma/cfd\_basilisk/rising\_bubble.c}. Он задает two-phase solver
Basilisk с поверхностным натяжением и параметрами rising-bubble benchmark.
Радиус меняется при компиляции с помощью макроса \texttt{RADIUS}, например:
\begin{center}
    \texttt{qcc -O2 -Wall -DRADIUS=0.20 rising\_bubble.c -o rising\_bubble -lm}
\end{center}
Каждый расчет сохраняет 151 снимок с шагом по времени около \(0.02\) на
интервале \(t\in[0,3]\). В исходных TSV-файлах Basilisk используется повернутая
полуобласть; скрипт \texttt{convert\_snapshots\_to\_h5.py} зеркалирует ее в
область \([-0.5,0.5]\times[0,2]\) и сохраняет данные в HDF5.

Важно отметить, что в статье \cite{Buhendwa2021} CFD-данные были получены
другим решателем, ALPACA. Поэтому локальные Basilisk-расчеты не должны
побитово совпадать с графиками из статьи; они используются как воспроизводимая
CFD-реализация того же benchmark-семейства. В Basilisk поле \(f\) соответствует
жидкой фазе, а в PINN используется \(\alpha\) для пузырька, поэтому при
конвертации берется \(\alpha=1-f\).

В HDF5-файлах хранятся массивы:
\[
    X,\quad Y,\quad time,\quad levelset,\quad density,\quad pressure,\quad
    velocityX,\quad velocityY.
\]
Размеры массивов для каждого радиуса:
\[
    |X|=256,\qquad |Y|=512,\qquad |time|=151.
\]
Файлы обучающей и тестовой выборок расположены в каталогах
\texttt{diploma/exp/train} и \texttt{diploma/exp/test}.

\begin{table}[h]
    \centering
    \begin{tabular}{c|c|c}
        \hline
        Набор & Файл & Радиус \(R\) \\
        \hline
        train & \texttt{rising\_bubble\_R018.h5} & 0.18 \\
        train & \texttt{rising\_bubble\_R022.h5} & 0.22 \\
        train & \texttt{rising\_bubble\_R025.h5} & 0.25 \\
        train & \texttt{rising\_bubble\_R028.h5} & 0.28 \\
        train & \texttt{rising\_bubble\_R032.h5} & 0.32 \\
        test  & \texttt{rising\_bubble\_R020.h5} & 0.20 \\
        test  & \texttt{rising\_bubble\_R030.h5} & 0.30 \\
        \hline
    \end{tabular}
    \caption{CFD-расчеты, используемые в параметрической задаче.}
    \label{tab:radii}
\end{table}

\subsection{Выбор обучающих точек}

Из полной временной истории выбираются 29 временных слоев:
\[
\begin{split}
    t\in\{&
    0,\ 0.02,\ 0.04,\ 0.30,\ 0.60,\ 0.70,\ldots,\ 3.00
    \}.
\end{split}
\]
Этот выбор повторяет идею статьи \cite{Buhendwa2021}: добавить плотные ранние
снимки \(t=0.02\) и \(t=0.04\), чтобы лучше восстановить начальную стадию,
и использовать более редкую сетку времени на оставшемся интервале.

Для каждого радиуса формируются следующие множества точек:
\[
    \mathcal{D}_\alpha,\qquad
    \mathcal{D}_f,\qquad
    \mathcal{D}_{bc}.
\]
Множество \(\mathcal{D}_\alpha\) используется для обучения поля объемной доли.
На каждом выбранном временном слое берется 500 точек около интерфейса и
400 точек в остальной области. Поэтому для одного радиуса
\[
    N_\alpha = 29(500+400)=26100.
\]

Множество \(\mathcal{D}_f\) используется для вычисления физических невязок.
В текущей реализации на каждом временном слое выбирается 400 точек
непосредственно около интерфейса, 2000 точек в окрестности интерфейса и
3000 точек в остальной области. Дополнительно к residual-точкам добавляются
граничные координаты, чтобы уравнения также контролировались на границах.
Такая стратегия является адаптацией схемы сгущения точек около интерфейса,
предложенной в \cite{Buhendwa2021}.

Для граничных условий используются:
\begin{itemize}
    \item точки на верхней границе для фиксации \(p_\infty^*=1\);
    \item пары точек на левой и правой границах для периодичности \(u\), \(v\)
    и \(p\);
    \item точки с нулевой скоростью на твердых границах, а также начальные
    боковые точки с нулевой скоростью.
\end{itemize}

\section{Параметрическая PINN-модель}

\subsection{Архитектура}

Нейронная сеть аппроксимирует отображение
\[
    \mathcal{N}_\theta:
    (x^*,y^*,t^*,R^*)\mapsto (u^*,v^*,p^*,\alpha).
\]
Входной параметр \(R^*\) позволяет использовать одну и ту же сеть для семейства
задач с разными начальными радиусами. Выходы сети задаются четырьмя независимыми
головами:
\[
    u = W_u h + b_u,
    \qquad
    v = W_v h + b_v,
\]
\[
    p = \exp\left(\operatorname{clip}(W_p h+b_p,-20,20)\right),
    \qquad
    \alpha=\operatorname{sigmoid}(W_\alpha h+b_\alpha).
\]
Экспонента гарантирует положительность давления, а сигмоида ограничивает
объемную долю диапазоном \(0\leq\alpha\leq1\).

В проекте реализованы несколько конфигураций. Основная сохраненная конфигурация
для параметрического эксперимента использует скрытые слои
\[
    32\rightarrow64\rightarrow128\rightarrow256\rightarrow128\rightarrow64\rightarrow32
\]
с функцией активации \(\tanh\). Также в коде предусмотрен более тяжелый
профиль \texttt{paper} с 8 скрытыми слоями по 350 нейронов.

\subsection{Функция потерь}

Полная функция потерь имеет вид
\begin{equation}
    \mathcal{L}
    =
    \mathcal{L}_\alpha
    +
    \mathcal{L}_{bc}
    +
    \lambda_m\mathcal{L}_m
    +
    \lambda_x\mathcal{L}_{mom,x}
    +
    \lambda_y\mathcal{L}_{mom,y}
    +
    \lambda_a\mathcal{L}_{\alpha,pde}.
    \label{eq:loss_total}
\end{equation}
Веса физических невязок совпадают с rising-bubble настройкой из статьи:
\[
    (\lambda_m,\lambda_x,\lambda_y,\lambda_a)=(1,10,10,1).
\]

Слагаемое по объемной доле:
\begin{equation}
    \mathcal{L}_\alpha
    =
    \frac{1}{N_\alpha}
    \sum_{i=1}^{N_\alpha}
    \left(
        \alpha_\theta(x_i^*,y_i^*,t_i^*,R_i^*)-\alpha_i
    \right)^2.
\end{equation}
Граничное слагаемое состоит из давления на верхней границе, периодичности
боковых границ и нулевой скорости в выбранных граничных/начальных точках:
\begin{equation}
    \mathcal{L}_{bc}
    =
    \mathcal{L}_{p,N}
    +
    \mathcal{L}_{EW}
    +
    \mathcal{L}_{u,v}.
\end{equation}
Для верхней границы
\begin{equation}
    \mathcal{L}_{p,N}
    =
    \frac{1}{N_N}
    \sum_{i=1}^{N_N}
    \left(p_\theta(\mathbf{z}_i^N)-p_\infty^*\right)^2,
    \qquad
    \mathbf{z}_i^N=(x_i^*,y_N^*,t_i^*,R_i^*).
\end{equation}
Для периодических границ
\begin{align}
    \mathcal{L}_{EW}
    =
    \frac{1}{N_{EW}}
    \sum_{i=1}^{N_{EW}}
    \big[
        &(u_\theta(\mathbf{z}_i^E)-u_\theta(\mathbf{z}_i^W))^2
        +(v_\theta(\mathbf{z}_i^E)-v_\theta(\mathbf{z}_i^W))^2
        \notag\\
        &+(p_\theta(\mathbf{z}_i^E)-p_\theta(\mathbf{z}_i^W))^2
    \big].
\end{align}
Для точек с заданной нулевой скоростью
\begin{equation}
    \mathcal{L}_{u,v}
    =
    \frac{1}{N_B}
    \sum_{i=1}^{N_B}
    \left[
        u_\theta(\mathbf{z}_i^B)^2
        +
        v_\theta(\mathbf{z}_i^B)^2
    \right].
\end{equation}

Физические слагаемые являются среднеквадратичными невязками:
\[
    \mathcal{L}_m
    =
    \frac{1}{N_f}\sum_{i=1}^{N_f} r_m(\mathbf{z}_i^f)^2,
    \qquad
    \mathcal{L}_{mom,x}
    =
    \frac{1}{N_f}\sum_{i=1}^{N_f} r_x(\mathbf{z}_i^f)^2,
\]
\[
    \mathcal{L}_{mom,y}
    =
    \frac{1}{N_f}\sum_{i=1}^{N_f} r_y(\mathbf{z}_i^f)^2,
    \qquad
    \mathcal{L}_{\alpha,pde}
    =
    \frac{1}{N_f}\sum_{i=1}^{N_f} r_\alpha(\mathbf{z}_i^f)^2.
\]
Все производные для \(r_m\), \(r_x\), \(r_y\) и \(r_\alpha\) вычисляются с
помощью automatic differentiation в PyTorch.

\subsection{Обучение}

Оптимизация выполняется алгоритмом Adam. Для основной параметрической серии
используется пять стадий обучения:
\[
\begin{array}{c|ccccc}
    \hline
    \text{Стадия} & 1 & 2 & 3 & 4 & 5 \\
    \hline
    \text{Эпохи} & 5000 & 5000 & 5000 & 5000 & 5000 \\
    \text{Learning rate} & 10^{-4} & 5\cdot10^{-5} & 10^{-5} &
    5\cdot10^{-6} & 10^{-6} \\
    \hline
\end{array}
\]
Данные перемешиваются в начале каждой эпохи. При конфигурации
\texttt{paper\_light} полный набор точек делится на 20 батчей.

\section{Вычислительный эксперимент}

\subsection{Схема проверки обобщения}

Обучение проводится только на радиусах \(R=0.18,0.22,0.25,0.28,0.32\).
После обучения модель применяется к двум HDF5-файлам из тестового набора:
\[
    R=0.20,\qquad R=0.30.
\]
Оба значения лежат внутри диапазона обучающих радиусов, но отсутствуют в
обучающей выборке. Поэтому эксперимент проверяет интерполяционную способность
модели по параметру \(R\).

Для сравнения с CFD-решением используются поля
\[
    u,\qquad v,\qquad p,\qquad \alpha.
\]
Давление сравнивается с учетом возможного сдвига калибровки, поскольку в
несжимаемом течении оно определено с точностью до константы. В коде визуализации
это реализуется выравниванием среднего давления на верхней границе.

\subsection{Метрики качества}

Для величины \(q\) на выбранной сетке можно использовать среднюю абсолютную
ошибку и относительные нормы:
\begin{align}
    MAE(q)
    &=
    \frac{1}{N}\sum_{i=1}^{N}|q_i^{\mathrm{PINN}}-q_i^{\mathrm{CFD}}|,\\
    L_1(q)
    &=
    \frac{\sum_{i=1}^{N}|q_i^{\mathrm{PINN}}-q_i^{\mathrm{CFD}}|}
         {\sum_{i=1}^{N}|q_i^{\mathrm{CFD}}|},\\
    L_2(q)
    &=
    \frac{\sqrt{\sum_{i=1}^{N}(q_i^{\mathrm{PINN}}-q_i^{\mathrm{CFD}})^2}}
         {\sqrt{\sum_{i=1}^{N}(q_i^{\mathrm{CFD}})^2}}.
\end{align}
В текущем скрипте \texttt{pinn/evaluate.py} дополнительно выводятся \(RMSE\),
\(MAE\) и относительный \(RMSE\).

\todo{Перед защитой добавить таблицу ошибок для \(R=0.20\) и \(R=0.30\):
\(MAE\), \(L_1\), \(L_2\) или \(RMSE\) для \(\alpha\), \(u\), \(v\), \(p\).
Числа нужно получить запуском \texttt{pinn/evaluate.py} на лучшем checkpoint.}

\begin{table}[h]
    \centering
    \begin{tabular}{c|cccc}
        \hline
        Радиус & \(\alpha\) & \(u\) & \(v\) & \(p\) \\
        \hline
        \(R=0.20\) & TODO & TODO & TODO & TODO \\
        \(R=0.30\) & TODO & TODO & TODO & TODO \\
        \hline
    \end{tabular}
    \caption{Ошибки PINN относительно CFD на радиусах, не использованных при обучении.}
    \label{tab:test_errors}
\end{table}

\subsection{Качественное сравнение}

Качественно модель должна воспроизводить три ключевых свойства CFD-решения:
\begin{enumerate}
    \item всплытие центра пузырька под действием силы Архимеда;
    \item деформацию интерфейса при сохранении связной формы пузырька;
    \item согласованные поля скорости и давления вокруг интерфейса.
\end{enumerate}

\begin{figure}[h]
    \centering
    \fbox{\parbox{0.85\textwidth}{
        TODO: вставить сравнение CFD и PINN для \(R=0.20\):
        строки \(\alpha,u,v,p\), столбцы CFD / PINN / ошибка,
        например для \(t=1.5\) и \(t=3.0\).
    }}
    \caption{Сравнение CFD и PINN для невидимого при обучении радиуса \(R=0.20\).}
    \label{fig:r020_comparison}
\end{figure}

\begin{figure}[h]
    \centering
    \fbox{\parbox{0.85\textwidth}{
        TODO: вставить сравнение CFD и PINN для \(R=0.30\):
        строки \(\alpha,u,v,p\), столбцы CFD / PINN / ошибка,
        например для \(t=1.5\) и \(t=3.0\).
    }}
    \caption{Сравнение CFD и PINN для невидимого при обучении радиуса \(R=0.30\).}
    \label{fig:r030_comparison}
\end{figure}

\begin{figure}[h]
    \centering
    \fbox{\parbox{0.85\textwidth}{
        TODO: вставить график истории обучения: total loss,
        \(\mathcal{L}_\alpha\), \(\mathcal{L}_{bc}\),
        \(\mathcal{L}_m\), \(\mathcal{L}_{mom,x}\),
        \(\mathcal{L}_{mom,y}\), \(\mathcal{L}_{\alpha,pde}\).
    }}
    \caption{История обучения параметрической PINN.}
    \label{fig:loss_history}
\end{figure}

\begin{figure}[h]
    \centering
    \fbox{\parbox{0.85\textwidth}{
        TODO: добавить график положения центра пузырька \(y_c(t)\)
        и, если возможно, скорости всплытия \(v_c(t)\) для CFD и PINN
        на тестовых радиусах.
    }}
    \caption{Траектория центра пузырька для тестовых радиусов.}
    \label{fig:center_of_mass}
\end{figure}

\subsection{Ожидаемая интерпретация результатов}

Если ошибки на \(R=0.20\) и \(R=0.30\) оказываются сопоставимыми с ошибками на
обучающих радиусах, это означает, что сеть выучила не только отдельные
траектории интерфейса, но и гладкую зависимость решения от параметра \(R\).
Наиболее чувствительными областями являются окрестность интерфейса и ранняя
стадия движения, где скорость близка к нулю, а относительные ошибки по скорости
могут быть искусственно завышены.

\todo{После получения численных результатов заменить этот абзац конкретным
выводом: какой тестовый радиус предсказывается лучше, где максимальная ошибка,
насколько хорошо совпадают форма пузырька и поле давления.}

\section{Заключение}

В работе построена параметрическая постановка PINN для задачи всплытия пузырька.
В отличие от исходной работы \cite{Buhendwa2021}, где сеть обучалась для одного
радиуса \(R=0.25\), здесь радиус включен во вход модели как параметр \(R^*\).
CFD-данные получены с помощью Basilisk для пяти обучающих и двух тестовых
радиусов. Математическая часть модели основана на безразмерных уравнениях
несжимаемого двухфазного течения с VOF-переменной, поверхностным натяжением и
гравитацией. Функция потерь объединяет данные о положении интерфейса,
граничные условия и физические невязки уравнений.

Ключевая особенность работы состоит в проверке способности PINN обобщать решение
по геометрическому параметру. Для завершения экспериментальной части необходимо
добавить численные ошибки и визуальные сравнения на радиусах \(R=0.20\) и
\(R=0.30\), которые не входят в обучающую выборку.

\begin{thebibliography}{9}

\bibitem{Buhendwa2021}
A. B. Buhendwa, S. Adami, N. A. Adams.
\textit{Inferring incompressible two-phase flow fields from the interface motion
using physics-informed neural networks}.
arXiv:2101.09833, 2021.
\url{https://arxiv.org/abs/2101.09833}.

\bibitem{Hysing2009}
S. Hysing, S. Turek, D. Kuzmin, N. Parolini, E. Burman, S. Ganesan, L. Tobiska.
\textit{Quantitative benchmark computations of two-dimensional bubble dynamics}.
International Journal for Numerical Methods in Fluids, 60(11), 1259--1288, 2009.

\bibitem{Hirt1981}
C. W. Hirt, B. D. Nichols.
\textit{Volume of fluid (VOF) method for the dynamics of free boundaries}.
Journal of Computational Physics, 39(1), 201--225, 1981.

\bibitem{Brackbill1992}
J. U. Brackbill, D. B. Kothe, C. Zemach.
\textit{A continuum method for modeling surface tension}.
Journal of Computational Physics, 100(2), 335--354, 1992.

\bibitem{Popinet2018}
S. Popinet.
\textit{Numerical models of surface tension}.
Annual Review of Fluid Mechanics, 50, 49--75, 2018.

\end{thebibliography}

\end{document}
